\TNchapter{Foreword: Why this book}
\label{ch:foreword}

\starttwocol

The agents are already on your desk.

Somewhere in your organisation this quarter, a slide deck is proposing one. It will probably be a customer-service agent, or a coding agent, or a procurement agent. It will have a budget number with at least six zeros and a launch date inside the year. It will have been shaped by a vendor who has done this before and by an internal team that mostly has not. And the question on the last slide will be a version of: \textit{do we go?}

We --- the authors of this book --- have watched a lot of teams answer that question. Some of the answers have been confident and correct. Most have been confident and underspecified. A few have been wrong in ways that turned into incidents nobody on the original team is still around to explain. The pattern that separates the first group from the other two is not technical sophistication. It is the discipline a sponsoring executive brings to the decision.

The discipline we are describing is what we call \textbf{ARIA}: a five-part way of thinking about the one question agents force on the enterprise --- \textit{how do we know we can trust this?} Alignment, hardening, benchmarking, evaluation, certification. Each is a body of practice. Each answers a different piece of the trust question. Together they are the difference between a competent agent program and an expensive accident.

This book is not for the people who build agents. There are good books for them, and the references at the end of every chapter point to the best of them. This book is for the executive who has to sponsor, approve, or kill an agent project --- and who needs the language to do it competently with her CFO, her CISO, her audit committee, and her board. We are writing for the person whose name will be on the decision.

That person, in this book, is Maya. She is a composite --- drawn from the dozens of VP-of-AI conversations that produced the material here --- but she is specific enough that we can write to her, and not at the field. Maya runs an AI portfolio at a mid-sized enterprise. She is technical enough to know what a large language model is and ambitious enough to be the one her CEO calls about agents. She is not a researcher, and she does not have to become one. What she needs is enough vocabulary to challenge her team without bluffing, enough framing to see the project's risks before her CISO does, and enough confidence to defend her decision afterwards.

The shape of the book follows the shape of the trust question. Part I (Alignment) asks how you tell an agent what you want. Part II (Hardening) asks how you keep it from doing damage. Part III (Benchmarking) asks how you know it is better than the alternative. Part IV (Evaluation) asks how you know it works for your job. Part V (Certification) asks how you prove you did the work. Every chapter ends with a short list of actions you can take with your team on Monday --- because the only book worth reading on a weekend is one you can act on the following week.

We have tried to be honest about what we do not know. The literature on agent safety is moving fast enough that some specifics in this book will be stale before its paperback edition. We have tried to write it so the framework outlasts the specifics. Where a 2026 technical detail does the heavy lifting, we have put it in a callout that a future editor can refresh without rewriting the surrounding argument. The pillars are designed to be stable; the names of the benchmarks they apply to are not, and that is fine.

A note on what this book is not. It is not a textbook on machine learning. It is not a vendor guide. It is not a regulatory cookbook, although it will help you read the cookbook your legal team hands you. It does not contain code. Where a topic has technical depth that we judge a sponsoring executive does not need, we have named the topic, glossed it in one sentence, and pointed at the literature. The cardinal rule of the prose is this: the body reads at the speed of an experienced colleague walking you through it. The technical depth lives in callouts and vignettes, where it can be skipped on a first read and returned to when a specific decision lands on your desk.

You should be able to read this book in a weekend. After that, the decision is yours.

\stoptwocol
